\DocumentMetadata{
  pdfversion=2.0,
  pdfstandard=ua-2,
  lang=en-US,
  tagging=on,
  tagging-setup={math/setup=mathml-SE,tabsorder=structure}
}
\documentclass[11pt,letterpaper]{article}
\usepackage[margin=0.68in,includefoot]{geometry}
\usepackage{newcomputermodern}
\usepackage{amsmath,amscd,mathtools}
\usepackage{tikz,adjustbox}
\providecommand{\square}{\mdlgwhtsquare}
\usepackage[hidelinks]{hyperref}
\hypersetup{
  pdftitle={Numerical Analysis PhD Exam, September 2002},
  pdfauthor={University of Florida Department of Mathematics},
  pdfsubject={Graduate mathematics examination},
  pdfdisplaydoctitle=true,
  bookmarksopen=true
}
\setcounter{secnumdepth}{0}
\setlength{\parindent}{0pt}
\setlength{\parskip}{0.28em}
\setlength{\emergencystretch}{3em}
\setlength{\leftmargini}{2.15em}
\setlength{\leftmarginii}{2.65em}
\setlength{\leftmarginiii}{3em}
\pagestyle{empty}
\raggedbottom
\allowdisplaybreaks
\NewTaggingSocketPlug{block/list/label}{examproblem}{%
  \tagstructbegin{tag=Lbl}\tagstructend%
  \tagstructbegin{tag=\LBody}%
  \tagstructbegin{tag=\ProblemHeadingTag,title-o={\ProblemHeadingText}}%
  \tagmcbegin{tag=\ProblemHeadingTag,actualtext={\ProblemHeadingText}}%
  #2%
  \tagmcend%
  \tagstructend%
}
\newenvironment{examproblems}[2]{%
  \def\ProblemHeadingTag{#2}%
  \AssignTaggingSocketPlug{block/list/label}{examproblem}%
  \begin{enumerate}%
  \setcounter{enumi}{#1}%
  \renewcommand{\labelenumi}{\textbf{\theenumi.}}%
  \setlength{\topsep}{0.35em}%
  \setlength{\partopsep}{0pt}%
  \setlength{\itemsep}{0.48em}%
  \setlength{\parsep}{0.18em}%
}{\end{enumerate}}
\newenvironment{examsubparts}{%
  \AssignTaggingSocketPlug{block/list/label}{default}%
  \begin{enumerate}%
  \setlength{\topsep}{0.18em}%
  \setlength{\partopsep}{0pt}%
  \setlength{\itemsep}{0.16em}%
  \setlength{\parsep}{0.1em}%
}{\end{enumerate}}
\newenvironment{examinstructions}{%
  \par\begingroup\itshape
}{%
  \par\endgroup\vspace{0.18em}
}
\makeatletter
\renewcommand\subsection{\@startsection{subsection}{2}{0pt}%
  {-1.25ex plus -.25ex minus -.1ex}%
  {0.55ex plus .1ex}%
  {\normalfont\large\bfseries}}
\renewcommand{\@maketitle}{%
  \newpage\null\vskip -1em%
  {\footnotesize\bfseries
    \href{https://www.ufl.edu/}{University of Florida}, \href{https://math.ufl.edu/}{Department of Mathematics}\par}%
  \vskip -0.5em%
  \tagstructbegin{tag=H1,title={Numerical Analysis PhD Exam, September 2002}}%
  \tagpdfparaOff%
  \tagmcbegin{tag=H1}%
    {\LARGE\bfseries\href{https://gma.math.ufl.edu/exams/numerical-analysis-phd/}{Numerical Analysis PhD Exam}, September 2002\par}%
  \tagmcend%
  \tagpdfparaOn%
  \tagstructend%
  \vskip 0.25em%
  \tagmcbegin{artifact}%
    \hrule height 1pt%
  \tagmcend%
  \vskip 1em%
}
\makeatother
\title{Numerical Analysis PhD Exam, September 2002}
\author{}
\date{}
\begin{document}
\maketitle
\thispagestyle{empty}
\begin{examinstructions}
Do any 8 of the following 10 problems.
\end{examinstructions}
\begin{examproblems}{0}{H2}
\def\ProblemHeadingText{Problem 1.}
\item
Consider a function \(f(x)\in C^2(\mathbb R)\) that satisfies the following properties:

\par
Using \(x_0=5/2\), do we know that Newton’s method will converge? If so, how many iterations are required to ensure \(10^{-4}\) accuracy?
\begin{examsubparts}
\item[\textnormal{(a)}]
There exists a unique root \(\alpha\in[2,3]\).
\item[\textnormal{(b)}]
For any real~\(x\), \(f'(x)\geq 3\) and \(0\leq f''(x)\leq 5\).
\end{examsubparts}
\def\ProblemHeadingText{Problem 2.}
\item
Let \(x_0,\ldots,x_n\) be distinct real points. Let
\[
P_n(x)=\sum_{j=0}^n c_j e^{jx}.
\]
 For given data \(y_0,\ldots,y_n\), show that there exists a unique choice of \(c_0,\ldots,c_n\) such that
\[
P_n(x_i)=y_i.
\]
 (Hint: Reduce to an ordinary interpolation problem)
\def\ProblemHeadingText{Problem 3.}
\item
Suppose you want to solve numerically the ode
\[
y'=f(x,y),\qquad y(0)=y_0.
\]
\begin{examsubparts}
\item[\textnormal{(a)}]
By applying Simpson’s rule to the integral in the equation
\[
y(x_{n+1})=y(x_{n-1})+\int_{x_{n-1}}^{x_{n+1}}f(s,y(s))\,ds,\qquad x_n=nh,\quad n=1,2,\ldots,
\]
 derive the multistep Simpson’s method for odes.
\item[\textnormal{(b)}]
Find the order of the local truncation error for this method.
\item[\textnormal{(c)}]
Determine if this method is stable or not.
\end{examsubparts}
\def\ProblemHeadingText{Problem 4.}
\item
This problem has two parts.
\begin{examsubparts}
\item[\textnormal{(a)}]
Find the linear least squares approximation to \(f(x)=x^3\) on the interval \([0,2]\) by optimizing the least squares error over~\(a\) and~\(b\) in \(r^*(x)=ax+b\).
\item[\textnormal{(b)}]
Find the first two orthonormal polynomials \(\phi_0(x),\phi_1(x)\) on \([0,2]\) (weight function \(w(x)=1\)). Use these to find the linear least squares approx to \(f(x)=x^3\), and check this result agrees with the previous problem.
\end{examsubparts}
\def\ProblemHeadingText{Problem 5.}
\item
This problem has two parts.
\begin{examsubparts}
\item[\textnormal{(a)}]
Let \(\mathbf A\in\mathbb C^{m\times n}\) with \(m\geq n\). Prove that \(\mathbf A^*\mathbf A\) is nonsingular if and only if \(\mathbf A\) has full rank.
\item[\textnormal{(b)}]
Let the \(\mathbf a_1,\mathbf a_2,\ldots,\mathbf a_\ell\in\mathbb C^m\) be linearly independent vectors. Form an \(m\times\ell\) matrix \(\mathbf A\) whose~\(i\)-th column is \(\mathbf a_i\). Prove that the matrix of the orthogonal projection onto span of \(\mathbf a_1,\mathbf a_2,\ldots,\mathbf a_\ell\) is
\[
\mathbf A(\mathbf A^*\mathbf A)^{-1}\mathbf A^*.
\]
\end{examsubparts}
\def\ProblemHeadingText{Problem 6.}
\item
Let~\(p\) and~\(q\) be positive real numbers such that \(1/p+1/q=1\), and let \(\|\cdot\|_p\) for \(1\leq p\leq\infty\) denote the norm on \(m\times n\) matrices induced by the \(\ell^p\)-norm on vectors in \(\mathbb C^m\) and \(\mathbb C^n\).
\begin{examsubparts}
\item[\textnormal{(a)}]
For any matrix, show that \(\|\mathbf A\|_p=\|\mathbf A^{\mathsf T}\|_q\). (Hint: use the Hölder inequality for vectors).
\item[\textnormal{(b)}]
Let \(\rho(\mathbf M)\) denote the spectral radius of any square matrix \(\mathbf M\). Prove that \(\rho(\mathbf M)\leq\|\mathbf M\|\) for any norm \(\|\cdot\|\) induced by a vector norm.
\item[\textnormal{(c)}]
Show that
\[
\|\mathbf A\|_2\leq\sqrt{\|\mathbf A\|_p\|\mathbf A^{\mathsf T}\|_q}.
\]
\end{examsubparts}
\def\ProblemHeadingText{Problem 7.}
\item
This problem has two parts.
\begin{examsubparts}
\item[\textnormal{(a)}]
Show that Jacobi’s iteration applied to \(\mathbf A\mathbf x=\mathbf b\) always converges when the matrix is row diagonally dominant.
\item[\textnormal{(b)}]
Show that the Gauss–Seidel iteration applied to \(\mathbf A\mathbf x=\mathbf b\) always converges when the matrix is row diagonally dominant.
\end{examsubparts}
\def\ProblemHeadingText{Problem 8.}
\item
If \(\mathbf A\) is a square matrix, then \(e^{\mathbf A}\) is the matrix obtained by forming the Taylor expansion of \(e^x\) and replacing~\(x\) by \(\mathbf A\). For any square matrix, show that \(\det e^{\mathbf A}\) is the product of the exponential of each eigenvalue of \(\mathbf A\).
\def\ProblemHeadingText{Problem 9.}
\item
If \(\mathbf A\) and \(\mathbf B\) are square invertible matrices and \(\mathbf u\) and \(\mathbf v\) are vectors with \(\mathbf A=\mathbf B-\mathbf u\mathbf v^{\mathsf T}\), obtain a formula for the scalar~\(\alpha\) in the following identity relating the inverses of \(\mathbf A\) and \(\mathbf B\):
\[
\mathbf A^{-1}=\mathbf B^{-1}+\alpha\mathbf B^{-1}\mathbf u\mathbf v^{\mathsf T}\mathbf B^{-1}.
\]
\def\ProblemHeadingText{Problem 10.}
\item
Let \(\mathbf A\) be an \(n\times n\) Hermitian positive definite matrix. Define
\[
\|\mathbf y\|_{\mathbf A}=(\mathbf y^*\mathbf A\mathbf y)^{1/2},\qquad \text{for all }\mathbf y\in\mathbb C^n.
\]
 Consider the following iterative method for solving \(\mathbf A\mathbf x=\mathbf b\):
\[
\mathbf x_{i+1}=\mathbf x_i+\alpha_i\mathbf r_i,\qquad \text{where }\mathbf r_i=\mathbf b-\mathbf A\mathbf x_i,\quad \text{and}\quad \alpha_i=\frac{\mathbf r_i^*\mathbf r_i}{\mathbf r_i^*\mathbf A\mathbf r_i}. \tag{1}
\]
 Let error at the~\(i\)-th step be \(\mathbf e_i=\mathbf x-\mathbf x_i\).
\begin{examsubparts}
\item[\textnormal{(a)}]
Prove that
\[
\|\mathbf e_{i+1}\|_{\mathbf A}=\inf_{\alpha\in\mathbb R}\|\mathbf e_i-\alpha\mathbf r_i\|_{\mathbf A}.
\]
\item[\textnormal{(b)}]
Use Problem (10a) to prove the following convergence rate estimate:
\[
\|\mathbf e_{i+1}\|_{\mathbf A}\leq\left(\frac{\kappa(\mathbf A)-1}{\kappa(\mathbf A)+1}\right)\|\mathbf e_i\|_{\mathbf A},
\]
 where \(\kappa(\mathbf A)\) denotes the spectral condition number of \(\mathbf A\).
\end{examsubparts}
\end{examproblems}
\end{document}
